\chapter{Phân tích và thiết kế}
\label{ch:thietke}

\begin{muctieu}
Chương này phân tích yêu cầu của Vị Việt, mô tả kiến trúc, dữ liệu, luồng
nhận diện và cấu hình model được tích hợp trong sản phẩm.
\end{muctieu}

\section{Phân tích bài toán}

Hệ thống cần giải quyết hai đường vào cùng hội tụ về một kho công thức. Với
tìm kiếm văn bản, người dùng nhập tên hoặc từ khóa; backend chuẩn hóa chuỗi
và trả danh sách món phù hợp. Với nhận diện ảnh, client gửi JPEG/PNG; model
sinh nhãn, backend chuẩn hóa nhãn đó rồi tìm bản ghi \code{dishes}. Đầu ra
cuối cùng là món chuẩn, ảnh minh họa, thời gian, khẩu phần, nguyên liệu, các
bước thực hiện và metadata nguồn.

Hai nhóm tác nhân được trình bày ở Hình~\ref{fig:usecase}. Khách sử dụng các
chức năng không cá nhân hóa. Người dùng đăng nhập có thêm yêu thích, lịch sử,
phản hồi và hồ sơ. Không có trang quản trị trong phạm vi đồ án; dữ liệu công
thức được cập nhật bằng catalog JSON và importer có kiểm tra.

\begin{figure}[htbp]
\centering
\resizebox{0.94\textwidth}{!}{\begin{tikzpicture}[x=1cm,y=1cm,
  actor/.style={rectangle, rounded corners=3pt, draw=daunavy,
    line width=.9pt, fill=lightnavy, minimum width=2.7cm,
    minimum height=10mm, align=center, font=\footnotesize, soft},
  usecase/.style={ellipse, draw=daumaroon, line width=.8pt, fill=white,
    minimum width=4.0cm, minimum height=10mm, align=center,
    font=\footnotesize}]
  \node[actor] (guest) at (-4.4,1.8) {Khách};
  \node[actor] (member) at (-4.4,-2.0) {Người dùng\\đăng nhập};

  \node[usecase] (browse) at (0,2.7) {Xem/tìm kiếm công thức};
  \node[usecase] (detail) at (0,1.2) {Xem chi tiết và nguồn};
  \node[usecase] (predict) at (0,-0.3) {Nhận diện món từ ảnh};
  \node[usecase] (auth) at (0,-1.8) {Đăng ký/đăng nhập};

  \node[usecase] (favorite) at (5.2,2.7) {Quản lý yêu thích};
  \node[usecase] (history) at (5.2,1.2) {Xem/xóa lịch sử};
  \node[usecase] (feedback) at (5.2,-0.3) {Gửi phản hồi nhãn đúng};
  \node[usecase] (profile) at (5.2,-1.8) {Sửa tên/đăng xuất};

  \node[draw=gray300, rounded corners=5pt, inner sep=8mm,
        fit=(browse)(profile), label={[font=\small\bfseries,
        text=daunavy]above:Hệ thống Vị Việt}] (system) {};

  \coordinate (guestbus) at (-2.55,1.8);
  \coordinate (memberbus) at (2.6,-3.0);
  \draw[assoc] (guest.east) -- (guestbus);
  \draw[assoc] (guestbus) |- (browse.west);
  \draw[assoc] (guestbus) |- (detail.west);
  \draw[assoc] (guestbus) |- (predict.west);
  \draw[assoc] (guestbus) |- (auth.west);

  \draw[assoc, -{Triangle[open,length=2.7mm,width=2.2mm]}]
    (member.north) -- (guest.south);
  \draw[assoc] (member.east) -- ++(1.1,0) |- (memberbus);
  \draw[assoc] (memberbus) |- (favorite.south);
  \draw[assoc] (memberbus) |- (history.south);
  \draw[assoc] (memberbus) |- (feedback.south);
  \draw[assoc] (memberbus) |- (profile.south);
\end{tikzpicture}
}
\caption{Sơ đồ Use Case tổng thể của Vị Việt.}
\label{fig:usecase}
\end{figure}

\section{Yêu cầu chức năng}

\begin{table}[htbp]
\centering
\caption{Các yêu cầu chức năng chính.}
\label{tab:functional-requirements}
\small
\begin{tabularx}{\textwidth}{>{\bfseries\color{daumaroon}}l X >{\centering\arraybackslash}p{2.2cm}}
\toprule
Mã & Mô tả & Tác nhân \\
\midrule
FR01 & Đăng ký, đăng nhập, refresh token và đăng xuất. & Người dùng \\
FR02 & Xem, tìm kiếm và lọc danh mục món. & Tất cả \\
FR03 & Xem chi tiết công thức, nguồn và giấy phép. & Tất cả \\
FR04 & Chụp/chọn ảnh, gửi nhận diện và xem Top~1. & Tất cả \\
FR05 & Mở trực tiếp công thức của món đã nhận diện. & Tất cả \\
FR06 & Thêm hoặc bỏ món yêu thích. & Đăng nhập \\
FR07 & Xem, xóa lịch sử và gửi phản hồi nhãn đúng. & Đăng nhập \\
FR08 & Đọc hồ sơ và thay đổi họ tên. & Đăng nhập \\
\bottomrule
\end{tabularx}
\end{table}

\section{Yêu cầu phi chức năng}

\begin{itemize}
  \item \textbf{Đúng dữ liệu:} không xuất bản công thức khi thiếu nguồn,
        nguyên liệu, bước nấu hoặc metadata ảnh.
  \item \textbf{Bảo mật:} mật khẩu băm Argon2; JWT có hạn dùng; refresh token
        xoay vòng và có thể thu hồi; secret production phải đủ mạnh.
  \item \textbf{Quyền riêng tư:} ảnh nhận diện không được lưu; lịch sử chỉ
        lưu nhãn, confidence, Top~3 và thời điểm.
  \item \textbf{Khả dụng:} giao diện tiếng Việt, vùng chạm đủ lớn, có trạng
        thái loading, lỗi, rỗng và fallback ảnh.
  \item \textbf{Tương thích:} Android Emulator dùng \code{10.0.2.2}; URL
        thiết bị thật cấu hình bằng \code{API\_BASE\_URL}.
  \item \textbf{Bảo trì:} database có migration; importer idempotent theo
        slug; model và dataset tách khỏi Git.
\end{itemize}

\section{Luồng nhận diện món ăn}

Hình~\ref{fig:activity} mô tả toàn bộ luồng từ lúc chọn ảnh đến khi mở công
thức. Ảnh sai định dạng hoặc không giải mã được bị chặn trước khi suy luận.
Nếu nhãn model không ánh xạ được tới database, client không tự chọn một món
khác. Đây là nguyên tắc quan trọng để tránh hiển thị công thức sai.

\begin{figure}[p]
\centering
\resizebox{!}{0.72\textheight}{\begin{tikzpicture}[node distance=7.5mm and 21mm]
  \node[startstop] (start) {Bắt đầu};
  \node[process, below=of start] (pick) {Chụp hoặc chọn ảnh};
  \node[decision, below=of pick] (valid) {Ảnh hợp lệ?};
  \node[process, right=of valid, minimum width=3.4cm] (error) {Thông báo lỗi\\và chọn lại};
  \node[process, below=of valid] (upload) {Gửi multipart /predict};
  \node[process, below=of upload] (infer) {ViT suy luận và lấy Top 3};
  \node[decision, below=of infer] (mapped) {Ánh xạ được món?};
  \node[process, right=of mapped, minimum width=3.4cm] (maperror) {Báo lỗi ánh xạ};
  \node[process, below=of mapped] (result) {Hiển thị Top 1\\và độ tin cậy};
  \node[decision, below=of result] (open) {Mở công thức?};
  \node[process, right=of open, minimum width=3.4cm] (retry) {Chụp ảnh khác};
  \node[process, below=of open] (recipe) {Mở đúng RecipeDetail};
  \node[startstop, below=of recipe] (end) {Kết thúc};

  \draw[flowarr] (start) -- (pick);
  \draw[flowarr] (pick) -- (valid);
  \draw[flowarr] (valid) -- node[above, font=\scriptsize]{Không} (error);
  \draw[flowarr] (error.north) |- (pick.east);
  \draw[flowarr] (valid) -- node[right, font=\scriptsize]{có} (upload);
  \draw[flowarr] (upload) -- (infer);
  \draw[flowarr] (infer) -- (mapped);
  \draw[flowarr] (mapped) -- node[above, font=\scriptsize]{Không} (maperror);
  \draw[flowarr] (maperror.east) -- ++(8mm,0) |- (end.east);
  \draw[flowarr] (mapped) -- node[right, font=\scriptsize]{có} (result);
  \draw[flowarr] (result) -- (open);
  \draw[flowarr] (open) -- node[above, font=\scriptsize]{Không} (retry);
  \draw[flowarr] (retry.east) -- ++(14mm,0) |- (pick.east);
  \draw[flowarr] (open) -- node[right, font=\scriptsize]{có} (recipe);
  \draw[flowarr] (recipe) -- (end);
\end{tikzpicture}
}
\caption{Activity Diagram của luồng nhận diện và mở công thức.}
\label{fig:activity}
\end{figure}
\clearpage

\section{Thiết kế kiến trúc hệ thống}

Kiến trúc ở Hình~\ref{fig:architecture} gồm ba vùng. Ứng dụng Flutter đảm
nhiệm giao diện, trạng thái phiên và gọi dịch vụ. FastAPI kiểm tra request,
xác thực, thực thi nghiệp vụ và chuẩn hóa response. Tầng tài nguyên gồm
SQLite, model ViT local, catalog công thức và 38 ảnh WebP.

\begin{figure}[htbp]
\centering
\resizebox{0.96\textwidth}{!}{\begin{tikzpicture}[node distance=12mm and 16mm]
  \node[objbox, minimum width=3.0cm, minimum height=12mm] (user)
    {Người dùng\\thiết bị Android};
  \node[clsbox, right=of user, minimum width=4.1cm, minimum height=12mm] (flutter)
    {Ứng dụng Flutter\\UI, AppState, services};
  \node[clsbox, right=of flutter, minimum width=4.1cm, minimum height=12mm] (api)
    {FastAPI backend\\router, validation, auth};

  \node[objbox, below=of flutter, xshift=-37mm,
        minimum width=4.0cm, minimum height=12mm] (media)
    {Dữ liệu công thức\\catalog + 38 ảnh WebP};
  \node[objbox, right=of media, minimum width=4.0cm,
        minimum height=12mm] (db)
    {SQLite/SQLAlchemy\\9 bảng nghiệp vụ};
  \node[objbox, right=of db, minimum width=4.0cm,
        minimum height=12mm] (model)
    {Model ViT-L/16 local\\processor + 38 nhãn};

  \node[draw=gray300, rounded corners=5pt, inner sep=5mm,
        fit=(flutter)(api), label={[font=\small\bfseries,
        text=daumaroon]above:Tầng ứng dụng và API}] {};
  \node[draw=gray300, rounded corners=5pt, inner sep=5mm,
        fit=(media)(model), label={[font=\small\bfseries,
        text=daunavy]below:Tầng dữ liệu và mô hình}] {};

  \draw[dataarr, <->] (user) -- node[above, font=\scriptsize]{thao tác} (flutter);
  \draw[flowarr, <->] (flutter) -- node[above, font=\scriptsize]{HTTP/JSON, multipart} (api);
  \draw[dataarr, <->] (api.south) -- ++(0,-6mm) -| (media.north);
  \draw[dataarr, <->] (api.south) -- ++(0,-9mm) -| (db.north);
  \draw[dataarr, <->] (api.south) -- ++(0,-12mm) -| (model.north);
  \draw[dataarr, dashed] (model) -- node[above, font=\scriptsize]{nhãn chuẩn hóa} (db);
\end{tikzpicture}
}
\caption{Kiến trúc tổng thể Flutter--FastAPI--SQLite--ViT.}
\label{fig:architecture}
\end{figure}

Model được nạp theo singleton và giữ trong bộ nhớ. Suy luận CPU được khóa ở
đoạn processor/model để tránh nhiều request đồng thời làm cạn RAM, trong khi
FastAPI đưa công việc sang worker thread để không chặn event loop. Media được
phục vụ qua \code{/media}; database chỉ lưu đường dẫn tương đối.

\clearpage
\section{Thiết kế API và trình tự xử lý}

Các router được chia theo auth, user, dish và activity. Biên API dùng JSON
camelCase; nội bộ Python và database dùng snake\_case. Request nhận diện là
multipart field \code{image}, chấp nhận JPEG/PNG tối đa 10~MB. Trình tự chi
tiết được minh họa ở Hình~\ref{fig:sequence}.

\begin{figure}[htbp]
\centering
\resizebox{0.96\textwidth}{!}{\begin{tikzpicture}[x=1cm,y=1cm, font=\scriptsize]
  \node[clsbox, minimum width=2.5cm] (client) at (0,0) {Flutter};
  \node[clsbox, minimum width=2.5cm] (api) at (4.4,0) {FastAPI};
  \node[clsbox, minimum width=2.5cm] (vit) at (8.8,0) {ViT-L/16};
  \node[clsbox, minimum width=2.5cm] (db) at (13.2,0) {SQLite};
  \foreach \x in {0,4.4,8.8,13.2} {
    \draw[dashed, gray500, line width=.55pt] (\x,-0.6) -- (\x,-7.25);
  }

  \draw[flowarr] (0,-1.0) -- node[above]{POST /predict + ảnh} (4.4,-1.0);
  \draw[flowarr] (4.4,-1.85) -- node[above]{processor(RGB 384)} (8.8,-1.85);
  \draw[dataarr] (8.8,-2.70) -- node[above]{logits, Softmax, Top 3} (4.4,-2.70);
  \draw[flowarr] (4.4,-3.55) -- node[above]{tìm normalized\_name} (13.2,-3.55);
  \draw[dataarr] (13.2,-4.40) -- node[above]{recognizedDish + recipe} (4.4,-4.40);
  \draw[flowarr] (4.4,-5.25) -- node[above]{nếu đăng nhập: lưu lịch sử} (13.2,-5.25);
  \draw[dataarr] (4.4,-6.10) -- node[above]{prediction + Top 3 + món} (0,-6.10);
  \draw[flowarr] (0,-6.72) -- ++(10mm,0) |- node[pos=.25, above]{mở công thức Top 1} (0,-7.18);
  \node[note, text width=5.8cm] at (7.8,-7.15)
    {Ảnh upload chỉ tồn tại trong bộ nhớ request, không ghi vào database hoặc ổ đĩa.};
\end{tikzpicture}
}
\caption{Sequence Diagram cho endpoint \code{POST /predict}.}
\label{fig:sequence}
\end{figure}

Response vẫn chứa Top~3 cho lịch sử và khả năng đánh giá, nhưng màn hình
chính chỉ hiển thị Top~1. Khi request có access token hợp lệ, backend tạo
\code{prediction\_history}; khi là khách, \code{predictionId} bằng null.

\clearpage
\section{Thiết kế cơ sở dữ liệu}

Database có chín bảng như Hình~\ref{fig:erd}. Bảng \code{dishes} là thực thể
trung tâm của kho công thức. Ba bảng con lưu nguyên liệu, bước nấu và nguồn
tham khảo theo thứ tự. \code{favorites} là bảng nối giữa user và dish. Lịch
sử và phản hồi lưu chuỗi nhãn thay vì khóa ngoại dish, bởi phản hồi có thể
đề cập nhãn model chưa ánh xạ được; vì vậy ERD không tạo quan hệ FK giả giữa
prediction và dishes.

\clearpage
\thispagestyle{dau}
\vspace*{\fill}
\begin{center}
\resizebox{\textwidth}{!}{\begin{tikzpicture}[node distance=10mm and 17mm,
  entity/.style={rectangle, rounded corners=2pt, draw=daunavy, fill=white,
    text width=4.45cm, minimum height=14mm, align=left, inner sep=5pt,
    font=\scriptsize, soft},
  rel/.style={-{Stealth[length=2mm]}, line width=.75pt, color=daumaroon},
  card/.style={font=\scriptsize\bfseries, fill=white, inner sep=1.5pt,
    text=daumaroon}]
  \node[entity] (users) {\textbf{users}\\\underline{id PK}, email UQ, full\_name, password\_hash, created\_at};
  \node[entity, below=of users] (sessions) {\textbf{refresh\_sessions}\\\underline{id PK}, jti UQ, user\_id FK, expires\_at, revoked};
  \node[entity, right=26mm of users] (favorites) {\textbf{favorites}\\\underline{id PK}, user\_id FK, dish\_id FK, created\_at\\UQ(user\_id,dish\_id)};
  \node[entity, right=26mm of favorites] (dishes) {\textbf{dishes}\\\underline{id PK}, slug UQ, name, normalized\_name, thời gian, khẩu phần, nguồn, ảnh, giấy phép, published};

  \node[entity, below=of favorites] (history) {\textbf{prediction\_history}\\\underline{id PK}, user\_id FK, predicted\_dish, confidence, top3\_json, selected\_dish, created\_at};
  \node[entity, below=of history] (feedback) {\textbf{prediction\_feedback}\\\underline{id PK}, history\_id FK UQ, user\_id FK, selected\_dish, created\_at};

  \node[entity, below=of dishes] (ingredients) {\textbf{recipe\_ingredients}\\\underline{id PK}, dish\_id FK, position, name, quantity, unit};
  \node[entity, below=of ingredients] (steps) {\textbf{recipe\_steps}\\\underline{id PK}, dish\_id FK, position, instruction};
  \node[entity, below=of steps] (refs) {\textbf{recipe\_references}\\\underline{id PK}, dish\_id FK, position, name, url, publisher, accessed\_at};

  \draw[rel] (users) -- node[card, left]{1:N} (sessions);
  \draw[rel] (users) -- node[card, above]{1:N} (favorites);
  \draw[rel] (dishes) -- node[card, above]{1:N} (favorites);
  \draw[rel] (users.east) -- ++(7mm,0) |- node[card, pos=.78, above]{1:N} (history.west);
  \draw[rel] (history) -- node[card, left]{1:0..1} (feedback);
  \draw[rel] (users.south) -- ++(0,-7mm) -| node[card, pos=.74, below]{1:N} (feedback.west);
  \draw[rel] (dishes) -- node[card, right]{1:N} (ingredients);
  \draw[rel] (dishes.east) -- ++(7mm,0) |- node[card, pos=.76, right]{1:N} (steps.east);
  \draw[rel] (dishes.east) -- ++(13mm,0) |- node[card, pos=.83, right]{1:N} (refs.east);
\end{tikzpicture}
}
\dauFigureCaption{ERD của chín bảng trong cơ sở dữ liệu Vị Việt.}{fig:erd}
\end{center}
\vspace*{\fill}
\clearpage

Foreign key dùng \code{ON DELETE CASCADE} ở các bản ghi phụ thuộc. Cặp
\code{user\_id,dish\_id} trong favorites là duy nhất; \code{history\_id}
trong feedback cũng duy nhất để mỗi lịch sử chỉ có tối đa một phản hồi.
Migration được version hóa bằng Alembic và chạy trước khi API nhận request.

\section{Thiết kế kho công thức và giấy phép}

Catalog được chia thành tám công thức dịch từ nguồn mở và ba mươi công thức
Vị Việt biên soạn. Nội dung dịch phải có URL phiên bản cố định, attribution,
ghi chú đã thay đổi và giấy phép CC BY-SA. Nội dung biên soạn cần tối thiểu
hai reference. Ảnh chỉ chấp nhận Public Domain, CC0, CC BY hoặc CC BY-SA và
phải có tác giả, URL nguồn, tên/URL giấy phép cùng SHA-256 trong manifest.

Importer chạy trong transaction, cập nhật theo slug và rollback toàn bộ nếu
một món không hợp lệ. Trước import thật, SQLite được sao lưu. Serializer chỉ
trả \code{hasRecipe=true} khi bản ghi thỏa toàn bộ điều kiện xuất bản.

\clearpage
\section{Thiết kế model nhận diện}

Pipeline ở Hình~\ref{fig:vit-pipeline} bắt đầu bằng kiểm tra MIME, dung lượng
và khả năng giải mã. Processor chính thức của model thực hiện resize/chuẩn
hóa; ứng dụng không tự crop theo một quy tắc khác. Cấu hình local cho biết
model có input 384$\times$384, patch 16, 24 encoder, hidden size 1024, 16
attention head và 38 lớp đầu ra.

\begin{figure}[htbp]
\centering
\resizebox{0.96\textwidth}{!}{\begin{tikzpicture}[node distance=10mm and 12mm]
  \node[io, minimum width=3.25cm] (image) {Ảnh JPEG/PNG\\tối đa 10 MB};
  \node[process, right=of image, minimum width=3.25cm] (decode) {Pillow kiểm tra\\và đổi RGB};
  \node[process, right=of decode, minimum width=3.25cm] (processor) {Image processor\\384$\times$384};
  \node[process, right=of processor, minimum width=3.25cm] (patch) {576 patch token\\+ CLS token};
  \node[process, below=of patch, minimum width=3.25cm] (encoder) {24 Transformer encoder\\hidden 1024, 16 head};
  \node[process, left=of encoder, minimum width=3.25cm] (softmax) {Linear 38 lớp\\+ Softmax};
  \node[process, left=of softmax, minimum width=3.25cm] (topk) {Top 3 backend\\Top 1 giao diện};
  \node[startstop, left=of topk, minimum width=3.25cm] (dish) {Tên món chuẩn\\và công thức};
  \draw[flowarr] (image) -- (decode);
  \draw[flowarr] (decode) -- (processor);
  \draw[flowarr] (processor) -- (patch);
  \draw[flowarr] (patch) -- (encoder);
  \draw[flowarr] (encoder) -- (softmax);
  \draw[flowarr] (softmax) -- (topk);
  \draw[flowarr] (topk) -- (dish);
\end{tikzpicture}
}
\caption{Pipeline suy luận của model ViT-L/16.}
\label{fig:vit-pipeline}
\end{figure}

Sau Softmax, backend trả Top~3 và cờ \code{accepted} dựa trên ngưỡng 0,60.
Tên nhãn được lấy trực tiếp từ \code{model.config.id2label}, chuẩn hóa dấu và
alias rồi truy vấn \code{dishes.normalized\_name}. Cách này tránh phụ thuộc
vào thứ tự lớp viết thủ công.

\refstepcounter{algorithm}
\label{alg:predict-recipe}
\addcontentsline{loa}{algorithm}{\protect\numberline{\thealgorithm}Nhận diện ảnh và ánh xạ công thức}
\begin{tcolorbox}[colback=gray100, colframe=daugold, boxrule=0.8pt,
  title={\textbf{Giải thuật~\thealgorithm. Nhận diện ảnh và ánh xạ công thức}},
  colbacktitle=daunavy, coltitle=white]
\small
\begin{tabularx}{\linewidth}{>{\bfseries\color{daumaroon}}r X}
Đầu vào & Ảnh $I$, model $M$, ngưỡng $\tau$, cơ sở dữ liệu $D$. \\
Đầu ra & Dự đoán Top~1, Top~3 và món được ánh xạ hoặc null. \\
1 & Nếu $I$ không phải JPEG/PNG hoặc vượt giới hạn thì trả lỗi đầu vào. \\
2 & $X \gets \operatorname{Processor}(\operatorname{RGB}(I))$. \\
3 & $P \gets \operatorname{Softmax}(M(X))$. \\
4 & $T \gets \operatorname{TopK}(P,3)$. \\
5 & $n \gets \operatorname{Normalize}(T[1].label)$. \\
6 & $d \gets \operatorname{FindDish}(D,n)$. \\
7 & Trả $T[1]$, $T$, $d$ và $(T[1].confidence \ge \tau)$. \\
\end{tabularx}
\end{tcolorbox}
